\section{Detector-systematics checks and image-dependent waveform phases}
\label{sec:systematics}

Repeated images provide useful consistency tests, but not a universal proof
that a frequency-dependent residual is astrophysical. The distinction is between
calibration that is stable between images, calibration that changes differently
in different instruments, and an epoch-dependent error common to the network.

\subsection{Apply the Morse operator before differencing}

In geometric optics, a macroimage contributes
\begin{equation}
 \widetilde h_i(f)=\sqrt{|\mu_i|}\,e^{2\pi ift_i}
 M_{n_i}(f)[1+\delta_i(f)]\widetilde h(f),\qquad
 M_n(f)=e^{-i\pi n\,\operatorname{sgn}f},
 \label{eq:sys_morse}
\end{equation}
using a convention in which a delay gives \(e^{2\pi ift_i}\).
The Morse index is \(n=0,1/2,1\) for a minimum, saddle, and maximum of
the arrival-time surface, respectively \cite{dai2017}. At nonzero frequencies,
\begin{center}
\begin{tabular}{lll}
Image type & Fourier operator & Effect on real strain \\
\hline
Minimum & \(1\) & unchanged \\
Saddle & \(-i\operatorname{sgn}f\) & Hilbert transform \\
Maximum & \(-1\) & overall sign reversal
\end{tabular}
\end{center}
Thus a saddle is not simply a sign-reversed waveform. Negative signed
magnification denotes image parity; the strain amplitude is
\(\sqrt{|\mu|}\), while its Morse phase is supplied separately. A maximum
has positive parity and nevertheless reverses the strain sign.
The relation \(M_n(-f)=M_n(f)^*\) preserves a real time-domain signal.
The zero-frequency geometric-optics limit is outside this calculation.

Before comparison, divide each image by its magnification and delay and apply
\(M_{n_i}^*(f)\). This removes the entire image operator, including for signals
with higher harmonics. Treating a saddle as a binary coalescence-phase shift is
not generally valid. An uncorrected Morse phase can otherwise produce a large
spurious residual. The fractional perturbation \(\delta_i\) remaining after
this correction has no universal sign; its amplitude and phase depend on the
perturbing structure and macroimage response \cite{oguri2022}.

The parity correction is not a free systematic-error reversal. A multiplicative
calibration error acts on the phase-shifted GW itself, so applying the inverse
Morse operator also rotates that error with the signal. After alignment, the
residual still contains $h[(\delta_i-\delta_j)+(c_{Ii}-c_{Ij})]$.
Image parity therefore helps identify and align lensed images; it does not
by itself distinguish diffraction from multiplicative calibration drift.

\subsection{Use independent detector observations of the same residual}

For a known sky direction, the network signal at each epoch can be written
\(\boldsymbol d_i=R_i[F_i\boldsymbol h]+\boldsymbol n_i\), where \(R_i\)
maps the two source polarizations into detector strains. At the scalar-lensing
order used here, \(F_i(f)\) multiplies both polarizations and is common to all
detectors observing that image. It therefore remains in the same two-polarization
signal subspace even if its frequency dependence is complicated. A noise-weighted
null projector satisfies
\begin{equation}
 Q_iR_i=0,\qquad Q_iR_i[F_i\boldsymbol h]=0.
\end{equation}
The null stream contains detector noise rather than the lensed GW. An
instrument-specific calibration error or glitch will generally leak into it.
This is the standard redundant-network consistency test
\cite{wenschutz2005,schutz2020}; it requires sufficient independent detector
responses and correct timing, sky location, and response modeling.

Image-to-image differencing supplies an additional check. For illustration,
suppose detector/polarization responses and the macro operators have been
corrected so that each scalar channel obeys
\begin{equation}
 x_{Ii}(f)=g_{Ii}(f)F_i(f)h(f),\qquad g_{Ii}=1+c_{Ii},
 \label{eq:sys_scalar}
\end{equation}
temporarily omitting additive noise. This is a model for response-corrected
channels, not a prescription for dividing raw strains whose antenna mixtures
change as Earth rotates. The image ratio in detector \(I\) is
\begin{equation}
 R_I^{ij}=\frac{x_{Ii}}{x_{Ij}}
 =\frac{F_i}{F_j}\frac{g_{Ii}}{g_{Ij}}.
\end{equation}
An arbitrary frequency-dependent calibration \emph{stable between epochs},
\(g_{Ii}=g_{Ij}\), cancels exactly. To first order,
\(R_I^{ij}-1=\delta_i-\delta_j+c_{Ii}-c_{Ij}\).
Taking the difference of this quantity between instruments cancels the lensing
signal and tests detector-dependent calibration evolution.

A division-free exact version for instruments \(I,J\) is the determinant
\begin{equation}
 \mathcal N_{IJ}^{ij}=x_{Ii}x_{Jj}-x_{Ij}x_{Ji}
 =F_iF_jh^2(g_{Ii}g_{Jj}-g_{Ij}g_{Ji}).
 \label{eq:sys_crossdifference}
\end{equation}
It vanishes for stable detector calibration even when the two images have
different diffraction. Its noise covariance must be propagated; in actual
analysis, fit the equivalent joint likelihood rather than divide noisy Fourier
bins. Agreement between detector-specific differential residuals, together with
a noise-consistent null stream, is strong evidence against an error confined to
one instrument.

\subsection{The remaining degeneracy is exact}

An epoch-dependent transfer \(a_i(f)\) common to every detector cannot be
separated from lensing by these tests. The transformation
\begin{equation}
 F_i(f)\longrightarrow a_i(f)F_i(f),\qquad
 g_{Ii}(f)\longrightarrow g_{Ii}(f)/a_i(f)
\end{equation}
leaves every measured strain unchanged. Such a common multiplicative error
passes the null stream and Eq.~\eqref{eq:sys_crossdifference}. External calibration
information and calibration priors remain necessary; a static common error
does cancel between images, but one that changes between their arrivals does not.

Physical diffraction templates also predict correlated amplitude and phase
structure and frequency coherence \cite{oguri2022}. Testing those predictions,
and checking that the same transfer acts on both polarizations, adds evidence.
Neither a smooth residual nor Hermitian symmetry \(F(-f)=F(f)^*\) is unique to
lensing: a physical detector transfer can satisfy both. There is no universal
minimum-phase amplitude--phase test for a multipath lens. Likewise,
\(R^{12}R^{23}R^{31}=1\) for three ratios formed from the same data is an
algebraic identity, not an independent lensing or calibration test.
